\chapter{Redox Reactions}
\def\shorttitle{Redox Reactions} 
\section[Background]{Background}

Understanding oxidation-reduction (redox) reactions that occur in natural 
and contaminated systems is central to developing accurate conceptual and 
quantitative models. Redox reactions involve the transfer of electrons 
between constituents and therefore, changes in the oxidation state of 
one or more elements.  Changes in oxidation state of an element can 
result in dramatic changes in its chemical and physical properties.  
For example, iron in the plus three oxidation 
state - Fe(III) - in surficial deposits renders 
a reddish hew to Mars whereas iron in the plus II oxidation 
state - Fe(II) - in surficial basalts on the Moon is, in no 
small part, responsible for the blackness of the lunar mares. 
Two important examples of redox reactions are the oxidation 
of aqueous Fe(II) to Fe(III) or organic carbon by dissolved oxygen:

\begin{equation}
4Fe^{+2} + O_{2,aq} + 10H_2O \rightarrow 4Fe(OH)_{3,s} + 8H^+  \\
\label{eq:fe2_oxidation}
\end{equation}

\begin{equation}
CH_2O  + O_{2,aq} + 10H_2O \rightarrow CO_{2,aq} + H_2O \\
\label{eq:om_reduction}
\end{equation}

Rather than representing formaldehyde, $CH_2O$ is often used to represent 
one sixth of a glucose molecule.

The classical approach to discussing redox reactions in aqueous geochemistry 
is to split the redox reactions into two hypothetical 
half reactions involving electrons. 
Using this approach, reaction \ref{eq:fe2_oxidation} and \ref{eq:om_reduction} 
would be broken up into the following half reactions:

\begin{equation}
Fe^{+2} + 3H_2O  \rightarrow Fe(OH)_{3,s} + 3H^+ + e^-  \\
\label{eq:fe_half_reac}
\end{equation}

\begin{equation}
CH_2O + H_2O \rightarrow CO_{2,aq} + 4H^+ + 4e^-  \\                    
\label{eq:om_half_reac}
\end{equation}

\begin{equation}
O_{2,aq} + 4H^+ + 4e^-  \rightarrow 2H_2O \\
\label{eq:ox_half_reac}
\end{equation}

Mass action expressions are then written for the half reactions.
For example, for the Fe(II)/Fe(III) half reaction, 
the mass action expression is:

\begin{equation}
K = a_{Fe(OH)_{3,s}} a_{H^+}^3 a_{e^-}/a_{Fe}^{+2} a_{H_2O}^3 \\
\label{eq:rdx5}
\end{equation}

The standard state of the $e^-$ is the same as other aqueous species,
1 mole per kilogram water.  
This approach comes from the field of electrochemistry, where the oxidation and 
reduction half reactions can be made to occur in separate place
connected, for example, 
through a meter that measures the electromotive force (or power source) 
and a salt bridge.
In dilute solutions the activity of water can be equated with its mole fraction,
which is essentially 1. If pure ferric hydroxide forms then it 
can be assumed to be in 
its standard state (mole fraction $Fe(OH)_{3,s}$
equals 1), 
so \ref{eq:rdx6} can be re-written:

\begin{equation}
K=a_{H^+}^3 a_{e^-}/a_{{Fe}^{+2}} \\    
\label{eq:rdx7}
\end{equation}
                                                         
Taking the base-10 logarithm, \ref{eq:rdx7} becomes: 
\begin{equation}
logK = 3loga_{H^+} + loga_{e^-} -loga_{{Fe}^{+2}}  \\                          
\label{eq:rdx8}
\end{equation}

Analogous to the definition of $pH$,
$loga_{e^-}$
is defined as $pe$.
This formalism is useful for graphically representing the 
relative thermodynamic stability of redox-reactive species 
in natural and contaminated systems 
(typically in their standard state).

Electrochemists have established logK values for 
half reactions based on the 
convention that the hydrogen gas -
hydrogen ion half reaction
\begin{equation}
H_{2,g}  =  2H^+ + 2e^- \\   
\label{eq:rdx9}
\end{equation} 
has an electromotive force and, therefore a logK 
value of 0 for $H_{2,g}$ and $H^+$ in their 
standard states (1 atmosphere and 1 mol/kgw, respectively).  
Using this convention, logK values for half reactions can be assembled 
in order of thermodynamic favorability, 
called a redox ladder, such as Table \ref{tab:rdx1}.  
Note that the logK values apply to all species in their standard state 
(1 mol/kgw for aqueous species, pure solid phase activity 
equals 1, and 1 atmosphere for gasses).  
Also, the logK value for half reactions where solid phases are participants 
depends on the thermodynamic stability of the solid phase.  
For example, reduction of ferrihydrite (fhy) is more favorable than 
reduction of goethite (based on the logK values for goethite that 
applies to the Cape Cod aquifer, Kent et al., 2008) because goethite 
is thermodynamically more stable than ferrihydrite.  
Plots of $pH$ against $pe$ constructed using reactions and logK values like 
those in Table \ref{tab:rdx1} are used to illustrate 
trends expected in redox processes 
with changes in pH. In the exercise that follows, we will explore these trends.

\begin{table}[h]
\topcaption{Half reactions and logK values for some redox-reactive species}
\centering
\footnotesize
%\begin{tabular*}{0.4\textwidth}{@{\extracolsep{\fill}}lll}
\begin{tabular}{ll}
\toprule
Half reactions, per electron	& logK \\
0.25O2,aq + e- + H+ = 0.5H2O	21.52 \\
0.2NO3- + e- + 1.2H+ = 0.1N2,aq + 0.6H2O	20.71 \\
0.5MnO2:H2Os,pyrolucite + e- + 2H+ = 0.5Mn+2 + 1.5H2O	20.69 \\
0.5HAsO4-2 + e- + 2H+   =  0.5As(OH)3,aq +  0.5H2O	19.80 \\
FeOOH + e- + 3H+ = Fe+2 + 2H2O  fhy	17.91 \\
FeOOH + e- + 3H+ = Fe+2 + 2H2O  Goethite (Kent et al. 2008)	13.52 \\
0.125SO4-2 + e- + 1.125H+ = 0.125HS- + 0.5H2O	4.21 \\
0.125CO3-2 + 1.25H+ + e- = 0.125CH4.aq + 0.375H2O	5.13 \\
\bottomrule
\label{tab:rdx1}
\end{tabular}
\end{table}	

\section[PHREEQC Exercise 6]{PHREEQC Exercise 6}
In this exercise we will examine the influence of chemical composition 
on redox processes and feedback between redox reactions and chemical composition.

Consider the half reactions listed in Table \ref{tab:rdx1}.  
Based on these half reaction fill in column two of \ref{tab:rdx2}. 
Enter the ratio of $H^+$ produced per $e-$ consumed.  
Use positive numbers where $H^+$ is produced by the consumption of 
electrons and negative numbers where $H^+$ is 
consumed by the consumption of electrons.  
This is, of course, not the conventional way of calculating stoichiometry, 
but what we are after is whether the redox process should 
decrease the pH (produce $H^+$), 
increase the pH (consume $H^+$), or have no impact on pH.  	

\begin{table}[h]
\topcaption{$H^+$ stoichiometry of redox reactions.}
\centering
\footnotesize
%\begin{tabular*}{0.4\textwidth}{@{\extracolsep{\fill}}lll}
\begin{tabular}{lll}
\toprule
Redox Couple	& $H^+$/$e^-$	& $H^+$/$C_6H_6$ \\
\midrule
$O_{2,aq}$/$H_2O$	         & & \\
$NO_3^-$/$N_{2,aq}$	       & & \\
$MnO_2·H_2O_{s}$/$Mn^{+2}$ & & \\
$H_2AsO_4^-$/$As(OH)_{3,aq}$ & & \\
$FeOOH_{s,fhy}$/$Fe^{+2}$ 	 & & \\
$FeOOH_{s,gte}$/$Fe^{+2}$ 	 & & \\
$SO_4^{-2}$/$CO_{2,aq}$ 		 & & \\
$CO_3^{-2}$/$CH_{4,aq}$      & & \\
\bottomrule
\label{tab:rdx2}
\end{tabular}
\end{table}	

Now, consider the half reaction where benzene is oxidized to $CO_{2,aq}$: 

\begin{equation}
C_6H_{6,aq} + 12H_2O   =  6CO_{2,aq} + 30e^- + 30H^+
\label{eq:rdx8}
\end{equation}

The $logK$ for this reaction is uncertain but we will assume that conversion 
of benzene to either $CO_2$ or methane is favorable.  
Combining this half reaction with each of the reactions in 
Table \ref{tab:rdx1} to eliminate $e-$ produces the following overall 
reactions for the oxidation of benzene by each of the oxidants.  
Using the reactions in Table \ref{tab:rdx3}, 
fill in column three of Table \ref{tab:rdx2}

\begin{table}[h]
\topcaption{Overall reactions for the oxidation of benzene by the }
\centering
\footnotesize
%\begin{tabular*}{0.4\textwidth}{@{\extracolsep{\fill}}lll}
\begin{tabular}{rll}
\toprule 
Overall reaction &  &\\
\midrule
$15O_{2,aq}$ + $2C_6H_{6,aq}$ + $6H_2O$    & = & $12CO_{2,aq}$ + $6H_2O$  = $12HCO_3^-$ + $12H^+$  \\
$6NO_3^-$ + $C_6H_{6,aq}$ + $6H^+$         & = & $3N_{2,aq}$ + $6CO_{2,aq}$ + $6H_2O$ \\
$15MnO_2·H_2O_s$ + $C_6H_{6,aq}$ + $30H^+$ & = & $15Mn^{+2}$ + $6CO_{2,aq}$ + $33H_2O$ \\
$15H_2AsO_4^-$ + $C_6H_{6,aq}$  + $15H^+$  & = & $15As(OH)_{3,aq}$ +  $3H_2O$ + $6CO_{2,aq}$ \\
$30FeOOH_{s,fhy} + C_6H_{6,aq} + 60H^+$    & = & $30Fe^{+2} + 48H_2O + 6CO_{2,aq}$ \\
$30FeOOH_{s,gte} + C_6H_{6,aq} + 60H^+$    & = & $30Fe^{+2} + 48H_2O + 6CO_{2,aq}$ \\
$15SO_4^{-2} + 4C_6H_{6,aq} + 15H^+$       & = & $24CO_{2,aq} + 15HS^- + 12H_2O$ \\
$4C_6H_{6,aq} + 18H_2O$                    & = & $15CH_{4.aq} + 9CO_{2,aq}$ \\
\bottomrule
\label{tab:rdx3}
\end{tabular}
\end{table}	

For which reactions are there qualitative changes in the $H^+$
stoichiometry (influence of pH) when the overall reaction 
is considered as compared to the half reaction ?

What quantitative differences are there in the
$H^+$ stoichiometry between Fe(III) reduction and sulfate 
reduction when the overall reaction is considered as compared to 
the half reaction ?  \\
Fe(III) reduction and methanogenesis ?

Calculate the number of umoles of $C_6H_6$ required to deplete 
each of the electron acceptors.  
Record the number of umoles of $C_6H_6$ to the tenths place and 
then tabulate the running total in Table \ref{tab:rdx4}.  
You will use the running total in the PHREEQC exercise that follows.

\begin{table}[h]
\topcaption{Hand-calculation table. EA refers electron acceptor.}
\centering
\footnotesize
%\begin{tabular*}{0.4\textwidth}{@{\extracolsep{\fill}}lll}
\begin{tabular}{llll}
\toprule 
EA	& Initial $\mu$moles & $\mu$moles $C_6H_6$ to deplete EA	& Total $\mu$moles $C_6H_6$ added \\
$O_2$	& 250		& & \\
$NO_3^-$	& 1500	& & \\
$MnO_{2,s}$	& 800	& & \\
$As(V)$	& 0.55 & & \\
$FeOOH_s$	& 8000 & & \\
$SO_4^{-2}$	& 260	& & \\
\bottomrule
\label{tab:rdx4}
\end{tabular}
\end{table}

Now we will use PHREEQC to conduct a redox titration 
(Abrams et al., 1996; Appelo and Postma, 2005).  
Small increments of benzene will be added to a suspension with 
one kilogram of water whose initial composition 
is listed in {\color{blue}redox\_titration.phrq}
% Table \ref{tab:rdx5}.
Based on the composition and thermodynamic data already described, 
in what order to you expect the electron acceptors 
to be completely reduced by benzene ?

We will use the PHREEQC file {\color{blue}redox\_titration.phrq}
%in Table \ref{tab:rdx5},
which adds benzene in a series of steps 
pre-determined to make the model data plot up well. The results of the redox titration are shown in 
Figure \ref{fig:redox_titration}.  \\
\\
In what order were the electron acceptors reduced ?  \\
\\
How does this order differ from the order predicted by the redox ladder (Table \ref{tab:rdx1}) ? \\
\\
Why ?

 

\begin{figure}
\centering
\includegraphics[width=1.0\textwidth]{redox_figure}
\caption[Results of the redox titration simulation.]{Results of the redox titration simulation.}
\label{fig:redox_titration}
\end{figure}








 

 


